\documentclass[14pt,aspectratio=1610]{beamer}
%\documentclass[14pt,aspectratio=1610,handout]{beamer}
\usepackage{csu-theme}
\usepackage{arrays}
\usepackage{booktabs} % For tables

\title{\Huge{} Searching \& Quadratic Sorting Algorithms}
\subtitle{CSCI 315: Data Structure Analysis\\Chapter 18}
\author{Sean T. Hayes, Ph.D.}
\institute{Charleston Southern University}

%% Comment out date to use default (today)
% \date{February 13, 2025}
\subject{Data-Strucutre Analysis}
\keywords{seach, sort, algorithm, performance analysis}

\setlength{\parskip}{0.5em}

\graphicspath{{../imgs/performance}}

\begin{document}

\begin{frame}
  \titlepage{}
\end{frame}
\section{Introduction}

\begin{frame}{Objectives}
After the next two lectures and labs, you will be able to:
\begin{itemize}
  \item
    Implement the sequential and binary search algorithms.
  \item
    Explain how these search algorithms perform.
  \item
    State lower bound on comparison-based search algorithms.
  \item
    Implement various sorting algorithms.
  \item
    Use the \textit{asymptotic} notation, Big-$O$, to analyze these algorithms' performance.
\end{itemize}
\end{frame}

\begin{frame}{With a Search Algorithm, you can:}
\begin{itemize}
  \item
    Determine whether a particular item is in an array or list.
  \item
    If the data is specially organized (for example, sorted), find the location in the list where a new item can be inserted.
  \item
    Find the location of an item to be deleted.
\end{itemize}
\end{frame}

\begin{frame}{Search Algorithm Terminology}
\begin{itemize}
  \item
    An item's \emph{key} -- a special member that uniquely identifies the item in the data set.

  \item
    A \emph{key comparison} -- comparing the key of the search item with the key of an item in the list.
    \begin{itemize}
      \item
        We can count the number of key comparisons as a measure of performance.
    \end{itemize}
  \item
    The key that we are searching for is called the \emph{target} or \emph{search item}.
\end{itemize}
\end{frame}

\section{Sequential Search}

\begin{frame}{Sequential Search (a.k.a. Linear Search)}
\begin{itemize}
  \item
    Same for both array-based and linked lists.
  \item
    Starts at the first element and examines each element until a match is found.
  \item
    Can be implemented iteratively (using a loop) or recursively.
\end{itemize}
\end{frame}

\begin{frame}{Key Comparisons of a Sequential Search}
\begin{itemize}[<+->]
  \item
    Remember, the speed of a computer does not affect the number of key comparisons required.
  \item
    Statements before and after the loop are executed only once, $O(1)$.
  \item
    Statements in the while loop repeated some number from 1 to $n$ times (where $n$ is the number of keys).
\end{itemize}
\end{frame}

\begin{frame}{Key Comparisons of a Sequential Search}
\begin{itemize}[<+->]
  \item
    If the search key is not in the list, there will be \(n\) comparisons.
  \item
    In the best case, the first key matches the search item with just one comparison.
  \item
    In the worst case, the target is not in the list, so there will be \(n\) comparisons.
  \item
    Average number of comparisons:
    \begin{equation*}
      \frac{1 + 2 + \cdots + n}{n} = \frac{1}{n} \cdot \frac{n (n + 1)}{2} = \frac{n + 1}{2} = \frac{1}{2}n + \frac{1}{2}
    \end{equation*}
\end{itemize}
\end{frame}

\section{Binary Search}

\begin{frame}{Binary Search (for Sorted Lists)}
\begin{itemize}[<+->]
  \item
    Binary search can be applied to \emph{sorted} arrays or lists.
  \item
    Uses the ``divide and conquer'' technique.
  \item
    Steps
    \begin{enumerate}
      \item
        First, compare the target to the middle element.
      \item
        If the search item is less than the middle element, restrict the search to the lower half of the list.
      \item
        Otherwise, restrict the search to the upper half of the list.
      \item
        Repeat until found.
    \end{enumerate}
\end{itemize}
\end{frame}


\begin{frame}{Binary Search}
\center
Find the location of \texttt{75} in this sorted array.

\begin{tikzpicture}
  \ArrayLabel{list}
  \alt<-3>{
    \ArrayList{4, 8, 19, 25, 34}
    \alt<-2>{\ArrayNode{39}}{\ArrayNode[MyPurple]{39}}
  }{
    \ArrayList[MyGreen]{4, 8, 19, 25, 34, 39}
  }


  \onslide<3-4>{\ArrayElLabel[val]{mid}}

  \alt<-4>{
    \ArrayList{45, 48, 66, 75, 89, 95}
  }{
    \alt<-5>{
      \ArrayList{45, 48}
      \ArrayNode[MyPurple]{66}
    }{
      \ArrayList[MyGreen]{45, 48, 66}
    }

    \onslide<5-6>{\ArrayElLabel[el66]{mid}}
    \alt<-8>{
      \ArrayNode{75}
    }{
      \ArrayNode[MyPurple]{75}
    }

    \alt<-6>{
      \ArrayNode{89}
      \ArrayNode{95}
    }{
      \alt<-7>{
        \ArrayNode[MyPurple]{89}
        \ArrayNode{95}
      }{
        \ArrayList[MyGreen]{89, 95}
      }
      \alt<-8>{
        \ArrayElLabel[el89]{mid}
      }{
        \ArrayElLabel[el75]{mid}
      }
    }

    %\ArrayNode{95}
  }
  \onslide<2-3>{\ArrayGroup{el4}{el95}{Search Space}}
  \onslide<4-5>{\ArrayGroup{el45}{el95}{Search Space}}
  \onslide<6-7>{\ArrayGroup{el75}{el95}{Search Space}}
  \onslide<8->{\ArrayGroup{el75}{el75}{Search Space}}
  \onslide<9>{}
\end{tikzpicture}
\end{frame}

\begin{frame}{The Performance of the Binary Search}
\begin{itemize}[<+->]
  \item
    Every iteration cuts the size of the search list in half.
  \item
    If list $L$ has $1024 = 2^{10}$ items, at most $11$ iterations are needed to find $x$.
  \item
    Every iteration makes two key comparisons.
    \begin{itemize}
      \item
        For $L$, binary search has at most $22$ key comparisons.
      \item
        Max number of comparisons = $2\log_2n+2$.
      \item
        The sequential search required $512$ key comparisons (average) to find if $x$ is in $L$.
    \end{itemize}
\end{itemize}
\end{frame}

\section{Big-O Notation}

\begin{frame}{Asymptotic Notation: Big-O Notation}
\begin{itemize}[<+->]
  \item
    After an algorithm is designed, it should be analyzed.
  \item
    There are often various ways to design a particular algorithm.
  \item
    Certain algorithms take very little computer time to execute.
  \item
    Others take a considerable amount of time.
\end{itemize}
\end{frame}

\begin{frame}{Asymptotic Notation: Big-O Notation}
\begin{table}[]
{
\rmfamily
\begin{tabular}{
>{}r r r r r }
\toprule
\multicolumn{1}{c}{$n$} & \multicolumn{1}{c}{$\log_2{n}$} & \multicolumn{1}{c}{$n \log_2{n}$} & \multicolumn{1}{c}{$n^2$} & \multicolumn{1}{c}{$2^n$} \\
\midrule
1               & 0     &   0                   & 1                & 2             \\
2               & 1     &   2                   & 4                & 4             \\
4               & 2     &   8                   & 16               & 16            \\
8               & 3     &  24                   & 64               & 256           \\
16              & 4     &  64                   & 256              & 65,536        \\
32              & 5     & 160                   & 1,024            & 4,294,967,296 \\
64              & 6     & 384                   & 4,096            & 1.84467e+19   \\
\bottomrule%
\end{tabular}
}
\end{table}
\end{frame}

\begin{frame}{Asymptotic Notation: Big-O Notation}
\center\vspace{-0.5em}%
\begin{table}[]
\caption{Time for $f(n)$ instructions if executing 1-billion instructions per second.}\vspace{-0.75em}
{\small%
\begin{tabular}{rlllll}
  \toprule%
$n$ & $n$ & $\log_2{n}$ & $n \log_2{n}$ & $n^2$ & $2^n$ \\
\midrule
$10$         & $0.01$\textmu{}s        & $0.003$\textmu{}s           & $0.033$\textmu{}s             & $0.1$\textmu{}s          & $1$\textmu{}s            \\
$20$         & $0.02$\textmu{}s        & $0.004$\textmu{}s           & $0.086$\textmu{}s             & $0.4$\textmu{}s          & $1$ms                \\
$30$         & $0.03$\textmu{}s        & $0.005$\textmu{}s           & $0.147$\textmu{}s             & $0.9$\textmu{}s          & $1$s                 \\
$40$         & $0.04$\textmu{}s        & $0.005$\textmu{}s           & $0.213$\textmu{}s             & $1.6$\textmu{}s          & $18.3$ min           \\
$50$         & $0.05$\textmu{}s        & $0.006$\textmu{}s           & $0.282$\textmu{}s             & $2.5$\textmu{}s          & $13$ days            \\
$100$        & $0.1$\textmu{}s         & $0.007$\textmu{}s           & $0.664$\textmu{}s             & $10$\textmu{}s           & $4\times10^{13}$ years       \\
$1000$       & $1$\textmu{}s           & $0.01$\textmu{}s            & $9.966$\textmu{}s             & $1$ms                &                    \\
$10000$      & $10$\textmu{}s          & $0.013$\textmu{}s           & $130$\textmu{}s               & $100$ms              &                    \\
$100000$     & $0.1$ms             & $0.017$\textmu{}s           & $1.67$ms                  & $10$s                &                    \\
$1000000$    & $1$ms               & $0.020$\textmu{}s           & $19.93$ms                 & $16.7$m              &                    \\
$10000000$   & $0.01$s             & $0.023$\textmu{}s           & $0.23$s                   & $1.16$ days          &                    \\
$100000000$  & $0.1$s              & $0.027$\textmu{}s           & $2.66$s                   & $115.7$ days\\
\bottomrule%
\end{tabular}}
\end{table}\vspace{-0.5em}
{\footnotesize *\textmu{}s is a microsecond or 10e-6 seconds.}
\end{frame}

\begin{frame}{Asymptotic Notation: Big-O Notation}
\begin{itemize}
  \item
    Let $f$ be a function of $n$.
  \item
    \emph{Asymptotic}: the study of the function $f$ as $n$ \\
    becomes larger and larger without bound.
  \item
    Let $f$ and $g$ be real-valued, non-negative functions.
  \item
    $f(n)$ is \emph{Big-O} of $g(n)$, written $f(n) = O(g(n))$ if there are constants $c$ and $n_0$ such that
     \begin{equation*}
      f(n) \le cg(n) \textsf{ for all } n \ge n_0
    \end{equation*}
\end{itemize}
\end{frame}

\begin{frame}{Growth Rate}
\begin{table}[]
{\rmfamily%
\begin{tabular}{rrr}
\toprule{}
  \(n\) & $g(n) = n^2$ & $f(n) = n^2 + 4n + 20$\\
\midrule{}
     10  &            100 &            160\\
     50  &          2,500 &          2,720\\
    100  &         10,000 &         10,420\\
  1,000  &      1,000,000 &     1,0004,020\\
 10,000  &    100,000,000 &    100,040,020\\
100,000  & 10,000,000,000 & 10,000,400,020\\
\bottomrule{}
\end{tabular}%
}%
\end{table}
\end{frame}

\begin{frame}{Growth Rate Functions}
\small
\begin{tabularx}{1\textwidth}{l X}
  \toprule{}
  \textbf{Function $g(n)$} & \textbf{Growth rate of $f(n)$}\\
  \midrule%
  $g(n) = 1$        & \emph{Constant} (is independent of the problem's size).\\
  $g(n) = \log_2n$  & \emph{Logarithmic} (increases increases slowly
                  as the problem size increases).\\
  $g(n) = n$        & \emph{Linear} (increases directly with the size
                  of the problem).\\
  $g(n) = n \log_2n$ & $n\log$ (increases more rapidly than a linear algorithm).\\
  $g(n) = n^2$      & \emph{Quadratic} (increases rapidly with the
                  size of the problem).\\
  $g(n) = n^3$      & \emph{Cubic} (increases more rapidly with the
                  size of the problem than the time requirement for a quadratic algorithm).\\
  $g(n) = 2^n$     & \emph{Exponential} (too rapid to be practical, as the size of the problem increases).\\
  \bottomrule{}
\end{tabularx}
\end{frame}

\begin{frame}{Comparing Sequential and Binary Search}

\begin{table}[]
\small
\begin{tabular}{r l l}
\toprule%\\[0.25ex]
\textbf{Algorithm} & \textbf{Successful Search} & \textbf{Unsuccessful Search}\\%[1.5ex]
%\hline{}\\[0.25ex]
\midrule
\textbf{Sequential Search} & $\frac{1}{2}n + \frac{1}{2} = O(n)$ & $n = O(n)$\\[1.5ex]
\textbf{Binary Search} & $2\log_2n - 3 = O(\log_2n)$ & $2\log_2n = O(log_2{n})$\\%[1.5ex]
\bottomrule
\end{tabular}
\caption{The number of comparisons for a list of length $n$.}
\end{table}
\end{frame}

\begin{frame}{Lower Bound on Comparison-Based Search}
\small
\emph{Comparison-based search algorithms}: Search a list by comparing \\the target element with list elements.
\pause{}

\emph{Theorem}: Let $L$ be a list of size $n > 1$. Suppose that the elements of $L$ are sorted. If $SRH(n)$ denotes the minimum number of comparisons needed, in the worst case, by using a comparison-based algorithm to recognize whether an element $x$ is in $L$, then $SRH(n) \ge \log_2(n + 1)$.
\pause{}

\emph{Corollary}: The binary search algorithm is an optimal worst-case algorithm for solving search problems by the comparison method.
From these results, it follows that if we want to design a search algorithm that is of an order less than $\log_2{n}$, then it cannot be comparison based.
\end{frame}

\begin{frame}{Sorting Algorithms}

\begin{itemize}
  \item
    To compare the performance of commonly used sorting algorithms, one must provide some analysis of these algorithms.
  \item
    These sorting algorithms can be applied to either array-based lists or linked lists.
\end{itemize}
\end{frame}

\section{Selection Sort}


\begin{frame}{Selection Sort}

\begin{itemize}
  \item<1->
    The \emph{Selection Sort Algorithm}: rearrange the list by selecting an element and moving it to its proper position.
  \item<2->
    Steps for a selection sort:
    \begin{itemize}
      \item
        Find the smallest element in the unsorted portion of the list
      \item
        Move it to the top of the unsorted portion by swapping it with the currently-first element.
      \item
        Start again with the rest of the list.
    \end{itemize}
  \item<3->
    Can also be applied to linked lists.
\end{itemize}
\end{frame}


\begin{frame}{Selection Sort}
\center
\begin{tikzpicture}[node of array/.append style={minimum width=26mm, inner sep=1pt}]
  \ArrayLabel{}

  % Index 0
  \alt<-3>{
    \ArrayNodeVertical{Sophia}
  }{
    \alt<-4>{
      \ArrayNodeVertical[MyPurple]{Aiden}
    }{
      \ArrayNodeVertical[MyGreen]{Aiden}
    }
  }

  % Index 1
  \alt<-7>{
    \ArrayNodeVertical{Jackson}
  }{
    \alt<-8>{
      \ArrayNodeVertical[MyPurple]{Ava}
    }{
      \ArrayNodeVertical[MyGreen]{Ava}
    }
  }

  % Index 2
  \alt<-11>{
    \ArrayNodeVertical{Olivia}
  }{
    \alt<-12>{
      \ArrayNodeVertical[MyPurple]{Emma}
    }{
      \ArrayNodeVertical[MyGreen]{Emma}
    }
  }
  \node[block] (elAt2) [fit=(val)]{};

  % Index 3
  \alt<-15>{
    \ArrayNodeVertical{Liam}
  }{
    \alt<-16>{
      \ArrayNodeVertical[MyPurple]{Emma}
    }{
      \ArrayNodeVertical[MyGreen]{Emma}
    }
  }


  % Index 4
  \alt<-9>{
    \ArrayNodeVertical{Emma}
  }{
    \alt<-11>{
      \ArrayNodeVertical[MyPurple]{Emma}
    }{
      \alt<-19>{
        \ArrayNodeVertical{Olivia}
      }{
        \alt<-20>{
          \ArrayNodeVertical[MyPurple]{Jackson}
        }{
          \ArrayNodeVertical[MyGreen]{Jackson}
        }
      }
    }
  }
  \onslide<10-11>{\ArrayElLabelVerticle[elEmma]{smallest}}
  \onslide<11>{\ArraySwapVerticle[50]{elOlivia}{elEmma}}
  \onslide<12>{\ArraySwapVerticle[50]{elEmma}{elOlivia}}

  % Index 5
  \alt<-13>{
    \ArrayNodeVertical{Emma}
  }{
    \alt<-15>{
      \ArrayNodeVertical[MyPurple]{Emma}
    }{
      \alt<-22>
      {
        \ArrayNodeVertical{Liam}
      }{
        \ArrayNodeVertical[MyGreen]{Liam}
      }
      
    }
  }

  % Index 6
  \alt<-5>{
    \ArrayNodeVertical{Ava}
  }{
    \alt<-7>{
      \ArrayNodeVertical[MyPurple]{Ava}
    }{
      \alt<-17>{
        \ArrayNodeVertical{Jackson}
      }{
        \alt<-19>{
          \ArrayNodeVertical[MyPurple]{Jackson}
        }{
          \alt<-24>{
            \ArrayNodeVertical{Olivia}
          }{
            \ArrayNodeVertical[MyGreen]{Olivia}
          }
        }
      }
    }
  }

  % Index 7
  \alt<1>{
    \ArrayNodeVertical{Aiden}
  }{
    \alt<-3>{
      \ArrayNodeVertical[MyPurple]{Aiden}
    }{
      \alt<-24>{
        \ArrayNodeVertical{Sophia}
      }{
        \ArrayNodeVertical[MyGreen]{Sophia}
      }
    }
  }

  \alt<-3>{
    \ArrayGroupVerticle{elSophia}{val}{Unsorted}
  }{
    \alt<-4>{
      \ArrayGroupVerticle{elAiden}{val}{Unsorted}
    }{
      \alt<-7>{
        \ArrayGroupVerticle{elJackson}{val}{Unsorted}
      }{
        \alt<-8>{
          \ArrayGroupVerticle{elAva}{val}{Unsorted}
        }{
          \alt<-11>{
            \ArrayGroupVerticle{elOlivia}{val}{Unsorted}
          }{
            \alt<-12>{
              \ArrayGroupVerticle{elAt2}{val}{Unsorted}
            }{
              \alt<-15>{
                \ArrayGroupVerticle{elLiam}{val}{Unsorted}
              }{
                \alt<-16>{
                  \ArrayGroupVerticle{elEmma}{val}{Unsorted}
                }{
                  \alt<-19>{
                    \ArrayGroupVerticle{elOlivia}{val}{Unsorted}
                  }{
                    \alt<-20>{
                      \ArrayGroupVerticle{elJackson}{val}{Unsorted}
                    }{
                      \alt<-22>{
                        \ArrayGroupVerticle{elLiam}{val}{Unsorted}
                      }{
                        \uncover<-24>{
                          \ArrayGroupVerticle{elOlivia}{val}{Unsorted}
                        }
                      }
                    }
                  }
                }
              }
            }
          }
        }
      }
    }
  }

  \onslide<2-3>{\ArrayElLabelVerticle[val]{smallest}}
  \onslide<3>{\ArraySwapVerticle[45]{elSophia}{val}}
  \onslide<4>{\ArraySwapVerticle[45]{elAiden}{val}}

  \onslide<6-7>{\ArrayElLabelVerticle[elAva]{smallest}}
  \onslide<7>{\ArraySwapVerticle[45]{elJackson}{elAva}}
  \onslide<8>{\ArraySwapVerticle[45]{elAva}{elJackson}}

  \onslide<14-15>{\ArrayElLabelVerticle[elEmma]{smallest}}
  \onslide<15>{\ArraySwapVerticle[45]{elLiam}{elEmma}}
  \onslide<16>{\ArraySwapVerticle[45]{elEmma}{elLiam}}

  \onslide<18-19>{\ArrayElLabelVerticle[elJackson]{smallest}}
  \onslide<19>{\ArraySwapVerticle[45]{elOlivia}{elJackson}}
  \onslide<20>{\ArraySwapVerticle[45]{elJackson}{elOlivia}}

  \onslide<22>{\ArrayElLabelVerticle[elLiam]{smallest}}
  \onslide<24>{\ArrayElLabelVerticle[elOlivia]{smallest}}
  \onslide<25>{}
\end{tikzpicture}
\end{frame}


\begin{frame}{Analysis: Selection Sort}
\begin{itemize}
  \item
    \texttt{swap}:\\
    Does three assignments; executed \(n - 1\) times.
    \begin{equation*}
      3(n - 1) = 3n - 3 = O(n)
    \end{equation*}
  \item<2->
    \texttt{minIndex}:
    \begin{itemize}
      \item
        For a list of length \(n\), \(n - 1\) key comparisons.
      \item
        Executed \(n - 1\) times (by \texttt{selectionSort})
      \item
        The number of key comparisons:
        \begin{equation*}
          (n-1)+(n-2)+\cdots+2+1=\frac{n(n-1)}{2}=\frac{1}{2} n^{2}-\frac{1}{2}n=O(n^{2})
        \end{equation*}
        \begin{equation*}
          O(n^2) + O(n) = O(n^2)
        \end{equation*}
  \end{itemize}
\end{itemize}
\end{frame}

\section{Bubble Sort}

\begin{frame}{Bubble Sort}
\begin{itemize}
  \item
    Suppose \lstinline{list[0]...list[n-1]} is a list of \(n\) elements, indexed \(0\) to \(n-1\).

  \item<2->
    \emph{Bubble sort algorithm}:
    \begin{itemize}
      \item
        In a series of $n-1$ iterations, compare successive elements, \lstinline{list[index]} and \lstinline{list[index + 1]}.
      \item
        If \lstinline{list[index]} is greater than \lstinline{list[index + 1]}, then swap them.
    \end{itemize}
\end{itemize}
\end{frame}

\begin{frame}{Bubble Sort}
\hspace{2cm}
\begin{tikzpicture}[node of array/.append style={minimum width=22mm, inner sep=1pt}]
  \ArrayLabel{}

  % Index 0
  \alt<-2>{%
    \ArrayNodeVertical{Emma}%
  }{%
    \alt<-15>{%
      \ArrayNodeVertical{Bella}%
    }{%
      \alt<-19>{%
        \ArrayNodeVertical{Aiden}%
      }{%
        \ArrayNodeVertical[MyGreen]{Aiden}%
      }%
    }%
  }%

  % Index 1
  \alt<-2>{%
    \ArrayNodeVertical{Bella}%
  }{%
    \alt<-11>{%
      \ArrayNodeVertical{Emma}%
    }{%
      \alt<-15>{%
        \ArrayNodeVertical{Aiden}%
      }{%
        \alt<-19>{%
          \ArrayNodeVertical{Bella}%
        }{%
          \ArrayNodeVertical[MyGreen]{Bella}%
        }%
      }%
    }%
  }%
  
  % Index 2
  \alt<-5>{%
    \ArrayNodeVertical{Liam}%
  }{%
    \alt<-11>{%
      \ArrayNodeVertical{Aiden}%
    }{%
      \alt<-17>{%
        \ArrayNodeVertical{Emma}%
      }{%
        \ArrayNodeVertical[MyGreen]{Emma}%
      }%
    }%
  }%

  % Index 3
  \alt<-5>{%
    \ArrayNodeVertical{Aiden}%
  }{%
    \alt<-7>{%
      \ArrayNodeVertical{Liam}%
    }{%
      \alt<-13>{%
        \ArrayNodeVertical{Jackson}%
      }{%
        \ArrayNodeVertical[MyGreen]{Jackson}%
      }%
    }%
  }%
  

  % Index 4
  \alt<-7>{%
    \ArrayNodeVertical{Jackson}%
  }{%
    \alt<-8>{%
      \ArrayNodeVertical{Liam}%
    }{%
      \ArrayNodeVertical[MyGreen]{Liam}%
    }%
  }%
  

  \onslide<2>{\ArraySwapVerticle<compare and swap>{elEmma}{elBella}}
  \onslide<3>{\ArraySwapVerticle<compare and swap>{elBella}{elEmma}}
  \onslide<4>{\ArraySwapVerticle<compare>{elEmma}{elLiam}}
  \onslide<5>{\ArraySwapVerticle<compare and swap>{elLiam}{elAiden}}
  \onslide<6>{\ArraySwapVerticle<compare and swap>{elAiden}{elLiam}}
  \onslide<7>{\ArraySwapVerticle<compare and swap>{elLiam}{elJackson}}
  \onslide<8>{\ArraySwapVerticle<compare and swap>{elJackson}{elLiam}}

  \onslide<10>{\ArraySwapVerticle<compare>{elBella}{elEmma}}
  \onslide<11>{\ArraySwapVerticle<compare and swap>{elEmma}{elAiden}}
  \onslide<12>{\ArraySwapVerticle<compare and swap>{elAiden}{elEmma}}
  \onslide<13>{\ArraySwapVerticle<compare>{elEmma}{elJackson}}

  \onslide<15>{\ArraySwapVerticle<compare and swap>{elBella}{elAiden}}
  \onslide<16>{\ArraySwapVerticle<compare and swap>{elAiden}{elBella}}
  \onslide<17>{\ArraySwapVerticle<compare>{elBella}{elEmma}}

  \onslide<19>{\ArraySwapVerticle<compare>{elAiden}{elBella}}

  \onslide<20>{}
\end{tikzpicture}
\end{frame}

\begin{frame}{Bubble Sort}
  \begin{itemize}
    \item
      The Bubble-Sort algorithm contains nested loops.
      \begin{itemize}
          \item
            The outer loop executes $n - 1$ times.
          \item
            For each iteration of the outer loop, the inner loop executes a certain number of times.
        \end{itemize}
    \item
      The total number of comparisons is:
      \begin{equation*}
        (n - 1) + (n - 2) + \cdots + 2 + 1 = \frac{n(n - 1)}{2} = \frac{1}{2}n^2 - \frac{1}{2}n = O(n^2)
      \end{equation*} 
    \item
      The number of assignments (worst case):
      \begin{equation*}
        3\frac{n(n - 1)}{2}= \frac{3}{2}n^2 - \frac{3}{2}n = O(n^2)
      \end{equation*}
      because there are 3 assignments per swap operator.
    \end{itemize}
\end{frame}

\section{Insertion Sort}

\begin{frame}{Insertion Sort}
The \emph{insertion sort algorithm} sorts the list by moving each element to its proper place in the sorted portion of the list.

\centerline{%
\begin{tikzpicture}[node of array/.append style={minimum width=18mm, inner sep=1pt}]
  \ArrayLabel{}
  \alt<1>{
    \ArrayList{10, 18, 25, 30, 23, 17, 45, 34}
  }{
    \alt<-2>{
      \ArrayNode[MyGreen]{10}
      \ArrayList{18, 25, 30, 23, 17, 45, 34}
    }{
      \alt<-3>{
        \ArrayList[MyGreen]{10, 18}
        \ArrayList{25, 30, 23, 17, 45, 34}
      }{
        \alt<-4>{
          \ArrayList[MyGreen]{10, 18, 25}
          \ArrayList{30, 23, 17, 45, 34}
        }{
          \ArrayList[MyGreen]{10, 18}
        }
      }
    }
  }
  \only<5-14>{
    \ArrayNode[MyGreen]{25}
  }
  \only<15->{
    \ArrayNode[MyGreen]{23}
  }

  \only<5->{
      \alt<-9>{
        \ArrayNode[MyGreen]{30}
        \only<5-9>{\ArrayGroup{el10}{el30}{Sorted}}
        \ArrayNode{23}
        \only<9>{\draw[link, ->] (el30) to[bend right=70] node[below] {copy} (el23);}
      }{
        \alt<-11>{
          \ArrayNode[MyGreen]{30}
          \onslide<10>{\draw[link, ->] (el30) to[bend right=70] node[below] {copy} (el23);}
          \onslide<11->{\draw[link, ->] (el25) to[bend right=70] node[below] {copy} (el30);}
          \onslide<5->{\ArrayGroup{el10}{el30}{Sorted}}
        }{
          \node[block] (first25) [fit=(val)]{};
          \ArrayNode[MyGreen]{25}
          \onslide<11-12>{\draw[link, ->] (first25) to[bend right=70] node[below] {copy} (el25);}
          \onslide<5-16>{\ArrayGroup{el10}{el25}{Sorted}}
        }
        \alt<-16>{
          \ArrayNode{30}
        }{
          \ArrayNode[MyGreen]{30}
        }
      }
      \ArrayList{17, 45, 34}
    }

  \onslide<21-22>{
    \draw[link, ->] (el30) to[bend right=70] node[below] {copy 1} (el17);
    \draw[link, ->] (el25) to[bend right=70] node[below] {copy 2} (el30);
    \draw[link, ->] (el23) to[bend right=70] node[below] {copy 3} (el25);
    \draw[link, ->] (el18) to[bend right=70] node[below] {copy 4} (el23);
  }

  \onslide<7->{
    \node[, name=prev] at (0.75, -2) {temp};
    \alt<7>{
      \node[node of array, anchor=west, name=val] at ([xshift=-1pt] prev.east) {};
    }{
      \alt<-19>{
        \node[node of array, anchor=west, name=val] at ([xshift=-1pt] prev.east) {23};
      }{
        \node[node of array, anchor=west, name=val] at ([xshift=-1pt] prev.east) {17};
      }
    }
  }

  \onslide<2>{\ArrayGroup{el10}{el10}{Sorted}}
  \onslide<3>{\ArrayGroup{el10}{el18}{Sorted}}
  \onslide<4>{\ArrayGroup{el10}{el25}{Sorted}}
  \onslide<17->{\ArrayGroup{el10}{el30}{Sorted}}
  \onslide<6>{\draw[link, <-, dotted] ([xshift=-2mm] el25.south east) to[bend right] node[below] {move} ([xshift=2mm] el23.south west);}
  \onslide<8>{\draw[link, <-] (val) to[bend right] node[above] {copy} (el23);}
  \only<14>{\draw[link, ->] (val) to[bend right] node[right] {copy} (first25);}
  \only<15>{\draw[link, ->] (val) to[bend right] node[right] {copy} (el23);}

  \onslide<18>{\draw[link, ->, dotted] (el17) to[bend left=50] node[below] {move} (el18);}
  \onslide<19-20>{\draw[link, <-] (val) to[bend right] node[above] {copy} (el17);}
  \onslide<22>{\draw[link, ->] (val) to[bend left] node[left] {copy} (el18);}
  
  \onslide<22>{}
\end{tikzpicture}%
}

\end{frame}

\begin{frame}{Insertion Sort Implementation for Arrays}
  \centering\Large
  Let's implement the Insertion Sort.
\end{frame}

% \begin{frame}[fragile]{Insertion Sort Implementation for Arrays}
% \begin{lstlisting}[basicstyle=\footnotesize\ttfamily]
% template <typename Type>
% void insertionSort(Type array[], unsigned int size)
% {
% 	for (auto firstUnsorted = 1u; firstUnsorted < size; ++firstUnsorted) {
% 		if (array[firstUnsorted] < array[firstUnsorted - 1]) {
% 			const auto TEMP = array[firstUnsorted];
% 			auto copyIndex = firstUnsorted;

% 			do {
% 				array[copyIndex] = array[copyIndex - 1];
% 				--copyIndex;
% 			} while(copyIndex > 0 && TEMP < array[copyIndex - 1]);

% 			array[copyIndex] = TEMP;
% 		}
% 	}
% }
% \end{lstlisting}
% \end{frame}

\begin{frame}{Insertion Sort}
  \begin{itemize}
    \item
      The for-loop executes $n - 1$ times.
    \item
      Best case (list is already sorted):\\
      Key comparisons: $n - 1 = O(n)$
    \item 
      Worst case: for each iteration, the if statement evaluates to true.\\
      Key comparisons: $1 + 2 + \cdots + (n - 1) = \frac{n(n - 1)}{2} = O(n^2)$
    \item 
      Average number of key comparisons and of item assignments:
      \begin{equation*}
        \frac{1}{4} n^2 + O(n) = O(n^2)
      \end{equation*}
    \end{itemize}
\end{frame}

\section{Lower Bound}

\begin{frame}{Sorting Algorithm Comparison}
\center{}
{\small Average Case Behavior of the Bubble Sort, Selection Sort, and Insertion Sort Algorithms for a List of Length $n$.}

{\small%
\begin{tabularx}{\textwidth}{Xcc}
\toprule{}
  \textbf{Algorithm} & \textbf{Comparisons} & \textbf{Swaps}\\
\midrule{}
Bubble Sort & $\frac{n(n-1)}{2} = O(n^2)$ & $\frac{n(n-1)}{4} = O(n^2)$ \\[1.3ex]
Selection Sort & $\frac{n(n-1)}{2} = O(n^2)$ & $3(n - 1) = O(n)$ \\[1.3ex]
Insertion Sort (Array Based) & $\frac{1}{4}n^2 + O(n) = O(n^2)$ & $\frac{1}{4}n^2 + O(n) = O(n^2)$ \\[1.3ex]
Insertion Sort (Linked-List Based) & $\frac{1}{4}n^2 + O(n) = O(n^2)$ & $n - 1 = O(n)$ \\
\bottomrule{}
\end{tabularx}%
}
\end{frame}

\begin{frame}{\large{Lower Bound on Comparison-Based Sort Algorithms}}
\center\Large
Can we do better than \(O(n^2)\) for sorting?\vspace{1em}

What is the best we can do?
\end{frame}

\section*{Lab 10}
\begin{frame}{\large{Lab 10: Selection Sort, Bubble Sort, and Binary Search}}
\centering\Large
  Let's take a look at Lab 10.
\end{frame}
\end{document}
